% Analysis section - interpretability and additional experiments

\subsection{What Do the Directions Encode?}
\label{sec:direction_analysis}

A key advantage of TDU over black-box unlearning methods is the interpretability of the learned directions. Each $\mathbf{f}_u^{(\ell)}$ is a unit vector in $\mathbb{R}^d$ that can be analyzed to understand what the model ``knows'' about the target fact.

\paragraph{Direction Similarity for Related Facts.}
We hypothesize that facts sharing semantic content should have similar encoding directions. To test this, we compute the cosine similarity between directions learned for related facts:

\begin{equation}
    \text{sim}(\mathbf{f}_u^{(i)}, \mathbf{f}_u^{(j)}) = \frac{\mathbf{f}_u^{(i)\top} \mathbf{f}_u^{(j)}}{\|\mathbf{f}_u^{(i)}\| \|\mathbf{f}_u^{(j)}\|}
\end{equation}

% TABLE: Direction similarity
\begin{table}[t]
\centering
\caption{Cosine similarity between directions for related vs. unrelated facts (Pythia-410M, layer 2).}
\label{tab:similarity}
\begin{tabular}{@{}lc@{}}
\toprule
\textbf{Fact Pair Type} & \textbf{Mean Cosine Similarity} \\
\midrule
Same subject, same relation & -- \\
Same subject, different relation & -- \\
Different subject, same relation & -- \\
Completely unrelated & -- \\
\bottomrule
\end{tabular}
\end{table}

Table~\ref{tab:similarity} shows that facts sharing semantic structure exhibit higher directional similarity. This suggests that TDU discovers genuine representational structure rather than arbitrary directions.

\paragraph{Token-Level Analysis.}
We analyze which tokens most strongly activate the learned directions by computing:
\begin{equation}
    a_t = \mathbf{f}_u^{\top} \mathbf{h}_t
\end{equation}
for each token position $t$ in a prompt. [Analysis to be added]

\subsection{Multi-Fact Unlearning}
\label{sec:multi_fact}

Practical applications require unlearning multiple facts simultaneously. We investigate two approaches:

\paragraph{Sequential Unlearning.}
Apply TDU iteratively for each fact. The rank-one updates compose:
\begin{equation}
    \mathbf{W}' = \prod_{i=1}^{n} \left(\mathbf{I} - \mathbf{f}_u^{(i)} \mathbf{f}_u^{(i)\top}\right) \mathbf{W}
\end{equation}

\paragraph{Joint Unlearning.}
Learn directions for multiple facts simultaneously by aggregating the forget set:
\begin{equation}
    \mathcal{D}_{\text{forget}} = \bigcup_{i=1}^{n} \mathcal{D}_{\text{forget}}^{(i)}
\end{equation}

% TABLE: Multi-fact scaling
\begin{table}[t]
\centering
\caption{Performance scaling with number of facts unlearned.}
\label{tab:multi_fact}
\begin{tabular}{@{}lccc@{}}
\toprule
\textbf{\# Facts} & \textbf{Mean ES} & \textbf{Mean SS} & \textbf{Retain $\Delta$} \\
\midrule
1 & -- & -- & -- \\
5 & -- & -- & -- \\
10 & -- & -- & -- \\
20 & -- & -- & -- \\
\bottomrule
\end{tabular}
\end{table}

Table~\ref{tab:multi_fact} shows that TDU scales reasonably to multiple facts, though we observe [interference effects / diminishing returns / etc. to be determined from experiments].

\subsection{Computational Efficiency}
\label{sec:efficiency}

% TABLE: Timing comparison
\begin{table}[t]
\centering
\caption{Wall-clock time for unlearning a single fact (GPU: NVIDIA A100).}
\label{tab:timing}
\begin{tabular}{@{}lccc@{}}
\toprule
\textbf{Method} & \textbf{GPT-2} & \textbf{Pythia-410M} & \textbf{Pythia-1B} \\
\midrule
Gradient Ascent (100 steps) & -- & -- & -- \\
ROME & -- & -- & -- \\
MEMIT & -- & -- & -- \\
TDU (ours) & -- & -- & -- \\
\bottomrule
\end{tabular}
\end{table}

Table~\ref{tab:timing} compares computational costs. TDU's efficiency stems from:
\begin{enumerate}
    \item No iterative optimization of model parameters
    \item Only forward/backward passes for direction learning
    \item Single rank-one weight update per layer
\end{enumerate}

\subsection{Limitations and Failure Cases}
\label{sec:limitations}

While TDU demonstrates strong performance, we identify several limitations:

\begin{itemize}
    \item \textbf{Complex facts:} Facts requiring multi-hop reasoning may not have a single encoding direction.
    \item \textbf{Entangled representations:} When facts share representational structure, removing one may affect related facts.
    \item \textbf{Verification:} While we can verify direction removal, confirming complete unlearning remains an open challenge.
\end{itemize}

We discuss these limitations and potential mitigations in Section~\ref{sec:conclusion}.
